\documentclass[12pt, letterpaper]{article}

\usepackage[utf8]{inputenc}
\usepackage[T1]{fontenc}
\usepackage[spanish]{babel}
\usepackage{geometry}
\usepackage{amsmath, amssymb, amsthm}
\usepackage{booktabs}
\usepackage{array}
\usepackage{microtype}
\usepackage{setspace}
\usepackage{parskip}
\usepackage{titlesec}
\usepackage{fancyhdr}
\usepackage{enumitem}
\usepackage{bm}
\usepackage{mathtools}
\usepackage{tcolorbox}
\usepackage{listings}
\usepackage{xcolor}
\usepackage{hyperref}

\tcbuselibrary{skins, breakable}

\geometry{top=2.8cm, bottom=3.0cm, left=3.2cm, right=3.2cm, headheight=14pt}
\setstretch{1.20}

% ---- Encabezado y pie de página ----
\pagestyle{fancy}
\fancyhf{}
\fancyhead[C]{\small\textsc{Econometría --- Gujarati --- Ejercicio 13.6}}
\fancyfoot[C]{\small\thepage}
\renewcommand{\headrulewidth}{0.4pt}
\renewcommand{\footrulewidth}{0pt}

% ---- Títulos de sección ----
\titleformat{\section}
  {\large\bfseries\scshape}{\thesection.}{0.6em}{}
  [\vspace{-4pt}\rule{\linewidth}{0.4pt}]
\titleformat{\subsection}{\normalsize\bfseries}{\thesubsection.}{0.5em}{}
\titleformat{\subsubsection}{\normalsize\itshape}{}{0em}{}

% ---- Caja resaltada ----
\newtcolorbox{clave}{
  colback=white, colframe=black,
  boxrule=0.5pt, arc=3pt,
  left=8pt, right=8pt, top=6pt, bottom=6pt,
  breakable
}

% ---- Estilo de código Stata ----
\definecolor{stataback}{RGB}{248,248,248}
\definecolor{stataframe}{RGB}{180,180,180}
\definecolor{statacomment}{RGB}{100,100,100}
\definecolor{statakeyword}{RGB}{0,70,140}

\lstdefinelanguage{Stata}{
  morekeywords={clear, set, gen, generate, quietly, regress, scalar,
                summarize, di, display, estimates, store, table,
                correlate, test, estat, log, close, seed, obs,
                noconstant, as, text, result, newline, if, replace,
                simulate, program, define, rclass, args, end,
                return, foreach, in, local, capture, drop},
  sensitive=false,
  morecomment=[l]{*},
  morecomment=[l]{//},
  morestring=[b]",
}

\lstset{
  language=Stata,
  backgroundcolor=\color{stataback},
  basicstyle=\ttfamily\footnotesize,
  keywordstyle=\color{statakeyword}\bfseries,
  commentstyle=\color{statacomment}\itshape,
  stringstyle=\color{black},
  frame=single,
  framesep=4pt,
  rulecolor=\color{stataframe},
  xleftmargin=8pt,
  xrightmargin=4pt,
  breaklines=true,
  breakatwhitespace=true,
  showstringspaces=false,
  tabsize=4,
  captionpos=b,
  numbers=none,
}

\newcommand{\Var}{\operatorname{Var}}
\newcommand{\E}{\operatorname{E}}

% ---- Entorno para salida de Stata ----
\lstdefinestyle{stataout}{
  language={},
  backgroundcolor=\color{white},
  basicstyle=\ttfamily\footnotesize,
  frame=single,
  framesep=4pt,
  rulecolor=\color{stataframe},
  xleftmargin=8pt,
  xrightmargin=4pt,
  breaklines=true,
  showstringspaces=false,
}

% ===========================================================================
\begin{document}
% ===========================================================================

\begin{center}
  {\LARGE\bfseries Ejercicio 13.6 --- Intercambio Sesgo-Varianza}\\[4pt]
  {\LARGE\bfseries mediante el Error Cuadrático Medio}\\[12pt]
  {\large\itshape Criterio del mínimo ECM para decidir entre un estimador}\\
  {\large\itshape sesgado con varianza pequeña y uno insesgado con varianza mayor}\\[14pt]
  \rule{0.55\linewidth}{0.6pt}\\[8pt]
  {\normalsize Ejercicio~13.6 --- \textit{Econometría}, Gujarati \& Porter, 5.a ed.\
  por Emanuel Quintana Silva}\\[4pt]
  {\small\textsc{Verificación empírica con datos simulados, $n = 200$, semilla 1306}}
\end{center}

\vspace{10pt}

%------------------------------------------------------------------------------
\section{Planteamiento del problema}
%------------------------------------------------------------------------------

Consulte las ecuaciones (13.3.4) y (13.3.5). Como se ve, $\hat{\alpha}_2$,
aunque sesgada, tiene una varianza menor que $\hat{\beta}_2$, que es
insesgada. ¿Cómo decidiría respecto de un intercambio de un sesgo por una
varianza pequeña?

\textbf{Sugerencia}: El ECM (error cuadrático medio) para los dos estimadores
se expresa como:

\begin{align}
  \mathrm{ECM}(\hat{\alpha}_2)
    &= \frac{\sigma^2}{\sum x_{2i}^2} + \beta_3^2\,b_{32}^2
    \quad \text{(varianza muestral + sesgo al cuadrado)} \label{eq:ecm_a}\\[6pt]
  \mathrm{ECM}(\hat{\beta}_2)
    &= \frac{\sigma^2}{\sum x_{2i}^2\,(1 - r_{23}^2)} \label{eq:ecm_b}
\end{align}

Este ejercicio es la síntesis del capítulo~13: en lugar de clasificar el error
de especificación como ``bueno'' o ``malo'', introduce el ECM como métrica
continua para decidir cuándo el sesgo es un precio razonable por reducir la
varianza.

\begin{center}
\renewcommand{\arraystretch}{1.3}
\small
\begin{tabular}{lccc}
\toprule
\textbf{Estimador} & \textbf{Sesgo} & \textbf{Varianza} & \textbf{Origen} \\
\midrule
$\hat{\alpha}_2$ (modelo corto) & $\beta_3\,b_{32}$ & $\sigma^2/\sum x_{2i}^2$ & Omite $X_3$ \\
$\hat{\beta}_2$ (modelo largo)  & $0$               & $\sigma^2/[\sum x_{2i}^2\,(1-r_{23}^2)]$ & Incluye $X_3$ \\
\bottomrule
\end{tabular}
\end{center}

%------------------------------------------------------------------------------
\section{El criterio del Error Cuadrático Medio}
%------------------------------------------------------------------------------

El ECM mide la exactitud total de un estimador al combinar su dispersión
aleatoria (varianza) y su error sistemático (sesgo):

\begin{equation}
  \mathrm{ECM}(\hat{\theta})
  = \Var(\hat{\theta}) + [\mathrm{Sesgo}(\hat{\theta})]^2
\end{equation}

La \textbf{regla de decisión} es: preferir el estimador con menor ECM, es
decir, aquel que minimiza el error total de predicción
(varianza$\,+\,$sesgo$^2$).

\subsection{ECM del estimador sesgado $\hat{\alpha}_2$ (modelo corto)}

\begin{equation}
  \mathrm{ECM}(\hat{\alpha}_2)
  = \underbrace{\frac{\sigma^2}{\sum x_{2i}^2}}_{\text{varianza}}
  + \underbrace{\beta_3^2\,b_{32}^2}_{\text{sesgo}^2}
\end{equation}

La varianza es la del estimador MCO simple, no afectada por multicolinealidad.
El sesgo cuadrático depende de cuán relevante sea $X_3$ ($\beta_3$) y de su
correlación con $X_2$ ($b_{32}$, pendiente de la regresión auxiliar de $X_3$
sobre $X_2$).

\subsection{ECM del estimador insesgado $\hat{\beta}_2$ (modelo largo)}

\begin{equation}
  \mathrm{ECM}(\hat{\beta}_2)
  = \Var(\hat{\beta}_2) + 0
  = \frac{\sigma^2}{\sum x_{2i}^2\,(1 - r_{23}^2)}
\end{equation}

El sesgo es cero. La varianza está inflada por el factor $1/(1-r_{23}^2)\geq 1$
(VIF), que crece con la correlación $r_{23}$ entre $X_2$ y $X_3$.

%------------------------------------------------------------------------------
\section{Condición de preferencia del estimador sesgado}
%------------------------------------------------------------------------------

Se prefiere $\hat{\alpha}_2$ si y solo si
$\mathrm{ECM}(\hat{\alpha}_2) < \mathrm{ECM}(\hat{\beta}_2)$. Sustituyendo
las expresiones~\eqref{eq:ecm_a} y~\eqref{eq:ecm_b} y restando
$\sigma^2/\sum x_{2i}^2$ en ambos lados:

\begin{equation}
  \beta_3^2\,b_{32}^2
  < \frac{\sigma^2}{\sum x_{2i}^2}
    \left(\frac{1}{1 - r_{23}^2} - 1\right)
  = \frac{\sigma^2}{\sum x_{2i}^2} \cdot \frac{r_{23}^2}{1 - r_{23}^2}
\end{equation}

\begin{clave}
  \textbf{Resultado}: el estimador sesgado tiene menor ECM cuando el cuadrado
  del sesgo es menor que la ganancia en varianza obtenida al omitir $X_3$:

  \begin{equation}
    \boxed{\beta_3^2\,b_{32}^2
      < \frac{\sigma^2\,r_{23}^2}{\sum x_{2i}^2\,(1 - r_{23}^2)}}
  \end{equation}
\end{clave}

%------------------------------------------------------------------------------
\section{Análisis del intercambio (tradeoff)}
%------------------------------------------------------------------------------

\subsection{Cuándo preferir el estimador sesgado $\hat{\alpha}_2$}

La omisión de $X_3$ es preferible cuando el sesgo introducido es pequeño en
relación con la ganancia en precisión:

\begin{enumerate}[leftmargin=1.8em]
  \item \textbf{$\beta_3 \approx 0$}: la variable omitida tiene efecto real
    débil sobre $Y$; el sesgo es despreciable.
  \item \textbf{$b_{32} \approx 0$}: $X_3$ está poco correlacionada con $X_2$;
    su omisión transfiere poco sesgo hacia $\hat{\alpha}_2$.
  \item \textbf{$r_{23}^2 \approx 1$}: alta multicolinealidad infla
    drásticamente la varianza de $\hat{\beta}_2$. Si el VIF es muy elevado,
    omitir $X_3$ puede reducir el ECM total a pesar del sesgo.
\end{enumerate}

\subsection{Cuándo preferir el estimador insesgado $\hat{\beta}_2$}

\begin{enumerate}[leftmargin=1.8em]
  \item \textbf{$\beta_3$ grande}: la variable omitida es teóricamente muy
    relevante.
  \item \textbf{$b_{32}$ grande}: $X_3$ tiene fuerte asociación con $X_2$;
    su omisión contamina severamente $\hat{\alpha}_2$.
  \item \textbf{Multicolinealidad moderada}: si $r_{23}^2$ no es cercano a~1,
    la inflación de varianza del modelo largo es manejable.
\end{enumerate}

\subsection{El papel del tamaño de muestra}

\begin{center}
\renewcommand{\arraystretch}{1.3}
\small
\begin{tabular}{p{2.2cm}p{3.2cm}p{3.2cm}p{4.8cm}}
\toprule
\textbf{Situación} & \textbf{Var($\hat{\beta}_2$)} &
\textbf{Sesgo($\hat{\alpha}_2$)} & \textbf{Recomendación} \\
\midrule
$n$ pequeño & Grande & Persistente & ECM puede favorecer sesgado \\
$n$ grande  & $\to 0$ & Persistente ($\neq 0$) &
  ECM dominado por sesgo; insesgado preferible \\
\bottomrule
\end{tabular}
\end{center}

En muestras grandes, las varianzas de estimadores consistentes tienden a cero,
pero el sesgo por omisión \textbf{no desaparece con} $n\to\infty$. El ECM de
$\hat{\alpha}_2$ converge al cuadrado del sesgo fijo, mientras el ECM de
$\hat{\beta}_2$ converge a cero.

\subsection{Resumen del intercambio}

\begin{center}
\renewcommand{\arraystretch}{1.3}
\small
\begin{tabular}{p{5.2cm}p{3.4cm}p{4.2cm}}
\toprule
\textbf{Condición} & \textbf{ECM favorece} & \textbf{Razón} \\
\midrule
$\beta_3$ pequeño, $b_{32}$ pequeño
  & $\hat{\alpha}_2$ (sesgado)  & Sesgo$^2 \ll$ ganancia en varianza \\
$r_{23}^2$ muy alto (multicolinealidad severa)
  & $\hat{\alpha}_2$ (sesgado)  & VIF infla varianza del insesgado \\
$\beta_3$ grande, $b_{32}$ grande
  & $\hat{\beta}_2$ (insesgado) & Sesgo$^2$ domina el ECM del corto \\
$n$ grande
  & $\hat{\beta}_2$ (insesgado) & Sesgo persiste; varianza $\to 0$ \\
\bottomrule
\end{tabular}
\end{center}

%------------------------------------------------------------------------------
\section{Verificación empírica con datos simulados}
%------------------------------------------------------------------------------

Se generan datos del proceso verdadero $Y_i = \beta_1 + \beta_2 X_{2i}
+ \beta_3 X_{3i} + u_i$ bajo distintos escenarios de $\beta_3$ y $r_{23}$
para comparar los ECM empíricos de $\hat{\alpha}_2$ y $\hat{\beta}_2$.

\subsection{Código Stata --- simulación principal ($\beta_3 = 1$, $r_{23} = 0.8$)}

\begin{lstlisting}[caption={Escenario base: sesgo moderado y correlación alta}]
clear
set seed 1306
set obs 200

* Parámetros verdaderos
local b1 = 2
local b2 = 3
local b3 = 1        // efecto real de X3
local r23 = 0.8     // correlación entre X2 y X3

* Generar variables
gen X2 = rnormal(5, 2)
gen X3 = `r23' * X2 + sqrt(1 - `r23'^2) * rnormal(0, 2)
gen u  = rnormal(0, 1)
gen Y  = `b1' + `b2' * X2 + `b3' * X3 + u

* ---- Regresión auxiliar: b32 ----
quietly regress X3 X2
scalar b32 = _b[X2]

* ---- Modelos corto y largo ----
quietly regress Y X2
scalar var_a2 = _se[X2]^2

quietly regress Y X2 X3
scalar var_b2 = _se[X2]^2

quietly correlate X2 X3
scalar r23_hat = r(rho)

* ---- ECM teórico ----
quietly regress Y X2 X3
scalar sigma2 = e(rmse)^2
quietly summarize X2
scalar sum_x2 = (r(N)-1) * r(Var)

scalar ecm_a2_teo = sigma2/sum_x2 + (`b3' * b32)^2
scalar ecm_b2_teo = sigma2/(sum_x2 * (1 - r23_hat^2))

* ---- Programa Monte Carlo ----
* Nota: if/else en bloques separados (evita r(198))
* Nota: cada local en su propia línea (';' no válido en .do)
capture program drop sim_ej136
program define sim_ej136, rclass
    args b3 r23
    clear
    set obs 200
    gen X2 = rnormal(5, 2)
    gen X3 = `r23' * X2 + sqrt(1 - `r23'^2) * rnormal(0, 2)
    gen u  = rnormal(0, 1)
    gen Y  = 2 + 3 * X2 + `b3' * X3 + u
    quietly regress Y X2
    return scalar a2 = _b[X2]
    quietly regress Y X2 X3
    return scalar b2 = _b[X2]
end

* ---- Simulación Monte Carlo: 1000 réplicas ----
simulate a2=r(a2) b2=r(b2), reps(1000) seed(1306): sim_ej136 `b3' `r23'

* ---- ECM empírico ----
quietly summarize a2
scalar ecm_a2_emp = r(Var) + (r(mean) - 3)^2
quietly summarize b2
scalar ecm_b2_emp = r(Var) + (r(mean) - 3)^2

* ---- Resultados ----
di as text _newline "=== Escenario: b3 = " `b3' ", r23 = " `r23' " ==="
di as text "b32 (regresion auxiliar):      " as result %8.4f b32
di as text "VIF = 1/(1-r23^2):             " as result %8.4f 1/(1 - r23_hat^2)
di as text _newline "--- ECM teorico ---"
di as text "ECM(alpha2) teorico:           " as result %8.6f ecm_a2_teo
di as text "ECM(beta2)  teorico:           " as result %8.6f ecm_b2_teo
di as text _newline "--- ECM empirico (MC, 1000 reps) ---"
di as text "ECM(alpha2) empirico:          " as result %8.6f ecm_a2_emp
di as text "ECM(beta2)  empirico:          " as result %8.6f ecm_b2_emp
di as text _newline "Estimador preferible (menor ECM):"
scalar _pref = (ecm_a2_emp < ecm_b2_emp)
if _pref == 1 {
    di as text "  --> alpha2 (sesgado) tiene menor ECM"
}
if _pref == 0 {
    di as text "  --> beta2 (insesgado) tiene menor ECM"
}
\end{lstlisting}

\subsection{Resultados --- escenario base}

\begin{center}
\renewcommand{\arraystretch}{1.3}
\small
\begin{tabular}{lrr}
\toprule
\textbf{Cantidad} & \textbf{Teórico} & \textbf{Empírico (MC, 1000 reps)} \\
\midrule
$b_{32}$ (regresión auxiliar)        & ---      & $0.7677$ \\
VIF $= 1/(1 - r_{23}^2)$             & ---      & $2.4451$ \\
\midrule
$\mathrm{ECM}(\hat{\alpha}_2)$        & $0.5907$ & $0.6412$ \\
$\mathrm{ECM}(\hat{\beta}_2)$         & $0.0035$ & $0.0035$ \\
\midrule
Estimador preferible & $\hat{\beta}_2$ (insesgado) & $\hat{\beta}_2$ (insesgado) \\
\bottomrule
\end{tabular}
\end{center}

Con $\beta_3 = 1$ y $r_{23} = 0.8$, el sesgo cuadrático
($\beta_3^2 b_{32}^2 \approx 0.59$) domina ampliamente la ganancia en
varianza ($\approx 0.002$); el modelo largo es claramente preferible.

\subsection{Código Stata --- comparación de cuatro escenarios}

\begin{lstlisting}[caption={Cuatro escenarios de tradeoff sesgo-varianza
  (requiere ejecutar el bloque anterior primero para definir sim\_ej136)}]
* ---- Comparación de cuatro escenarios ----
* Nota: sim_ej136 debe estar definido; se redefine aquí para asegurar
* que el bloque sea autónomo cuando se ejecuta desde un .do separado.
capture program drop sim_ej136
program define sim_ej136, rclass
    args b3 r23
    clear
    set obs 200
    gen X2 = rnormal(5, 2)
    gen X3 = `r23' * X2 + sqrt(1 - `r23'^2) * rnormal(0, 2)
    gen u  = rnormal(0, 1)
    gen Y  = 2 + 3 * X2 + `b3' * X3 + u
    quietly regress Y X2
    return scalar a2 = _b[X2]
    quietly regress Y X2 X3
    return scalar b2 = _b[X2]
end

* Escenario A: sesgo pequeño  (b3=0.1), correlación moderada (r23=0.50)
* Escenario B: sesgo grande   (b3=2.0), correlación moderada (r23=0.50)
* Escenario C: sesgo moderado (b3=1.0), correlación alta     (r23=0.95)
* Escenario D: sesgo moderado (b3=1.0), correlación baja     (r23=0.10)
* Nota: cada local en su propia línea (';' no es separador en .do)
local b3_A  = 0.1
local r23_A = 0.50
local b3_B  = 2.0
local r23_B = 0.50
local b3_C  = 1.0
local r23_C = 0.95
local b3_D  = 1.0
local r23_D = 0.10

foreach scen in A B C D {
    local b3  = `b3_`scen''
    local r23 = `r23_`scen''

    simulate a2=r(a2) b2=r(b2), reps(1000) seed(1306): ///
        sim_ej136 `b3' `r23'

    quietly summarize a2
    local ecm_a = r(Var) + (r(mean) - 3)^2
    quietly summarize b2
    local ecm_b = r(Var) + (r(mean) - 3)^2

    scalar _pref_`scen' = (`ecm_a' < `ecm_b')

    di as text "Escenario `scen': b3=`b3', r23=`r23' | " ///
       "ECM_a=" as result %7.5f `ecm_a' as text ///
       "  ECM_b=" as result %7.5f `ecm_b'
}

* ---- Preferencias (if en bloques separados, evita r(198)) ----
di as text _newline "=== Estimador preferible por escenario ==="
foreach scen in A B C D {
    if _pref_`scen' == 1 {
        di as text "  Escenario `scen': alpha2 (sesgado)"
    }
    if _pref_`scen' == 0 {
        di as text "  Escenario `scen': beta2  (insesgado)"
    }
}
\end{lstlisting}

\subsection{Interpretación de los cuatro escenarios}

\paragraph{Escenario A --- sesgo pequeño ($\beta_3 = 0.1$), correlación
moderada ($r_{23} = 0.5$).}
El sesgo cuadrático es ínfimo ($\beta_3^2 b_{32}^2 \approx 0.0025$) y el VIF
es solo $1.33$. El ECM del modelo corto está dominado por su varianza base,
prácticamente igual a la del modelo largo. En este caso
$\mathrm{ECM}(\hat{\alpha}_2) < \mathrm{ECM}(\hat{\beta}_2)$: el estimador
sesgado es preferible porque apenas hay sesgo que pagar.

\paragraph{Escenario B --- sesgo grande ($\beta_3 = 2.0$), correlación
moderada ($r_{23} = 0.5$).}
El sesgo cuadrático domina el ECM del modelo corto
($\beta_3^2 b_{32}^2 \approx 1.0$). Aunque el VIF es solo $1.33$, el cuadrado
del sesgo supera con creces la ganancia en varianza:
$\mathrm{ECM}(\hat{\alpha}_2) \gg \mathrm{ECM}(\hat{\beta}_2)$ y se prefiere
el estimador insesgado.

\paragraph{Escenario C --- sesgo moderado ($\beta_3 = 1.0$), correlación alta
($r_{23} = 0.95$).}
El VIF $= 1/(1 - 0.95^2) \approx 10.3$ infla severamente la varianza de
$\hat{\beta}_2$. A pesar del sesgo moderado, la reducción de varianza al omitir
$X_3$ puede dominar. El resultado Monte Carlo resuelve la ambigüedad: si
$b_{32}$ es grande (probable dado $r_{23} = 0.95$), el sesgo también será
grande y el modelo largo seguirá siendo preferible.

\paragraph{Escenario D --- sesgo moderado ($\beta_3 = 1.0$), correlación baja
($r_{23} = 0.1$).}
El VIF $\approx 1.01$ es prácticamente 1; incluir $X_3$ casi no infla la
varianza. El sesgo cuadrático domina y se prefiere sin duda el estimador
insesgado $\hat{\beta}_2$.

%------------------------------------------------------------------------------
\section{Conexión con los ejercicios anteriores del Capítulo 13}
%------------------------------------------------------------------------------

\begin{center}
\renewcommand{\arraystretch}{1.4}
\small
\begin{tabular}{p{1.0cm}p{4.0cm}p{2.2cm}p{2.4cm}p{3.0cm}}
\toprule
\textbf{Ej.} & \textbf{Error de especificación} & \textbf{Insesgadez} &
\textbf{Eficiencia} & \textbf{Criterio de elección} \\
\midrule
13.2 & Intercepto innecesario incluido     & Preservada  & Reducida           & Parsimonia con respaldo teórico \\
13.3 & Intercepto relevante omitido        & Perdida     & N/A                & Incluir siempre \\
13.4 & Variable irrelevante $X_3$ incluida & Preservada  & Reducida (VIF)     & Excluir con respaldo teórico \\
13.5a & Insumo $L_2$ relevante omitido    & Perdida     & N/A                & Incluir con respaldo teórico \\
13.5b & Insumo $L_2$ irrelevante omitido  & Preservada  & Mayor eficiencia   & Excluir; modelo corto correcto \\
\textbf{13.6} & \textbf{Tradeoff sesgo-varianza (ECM)} & \textbf{Depende} &
  \textbf{Mínimo ECM} & \textbf{Comparar ECM; teoría + $n$} \\
\bottomrule
\end{tabular}
\end{center}

%------------------------------------------------------------------------------
\section{Conclusión}
%------------------------------------------------------------------------------

\begin{clave}
  La decisión entre $\hat{\alpha}_2$ (sesgado, varianza menor) y
  $\hat{\beta}_2$ (insesgado, varianza mayor) no es mecánica. El criterio del
  \textbf{mínimo ECM} guía al investigador de la siguiente manera:

  \begin{equation*}
    \text{Preferir } \hat{\alpha}_2 \iff
    \beta_3^2\,b_{32}^2 <
    \frac{\sigma^2\,r_{23}^2}{\sum x_{2i}^2\,(1-r_{23}^2)}
  \end{equation*}

  En la práctica:
  \begin{enumerate}[leftmargin=1.4em]
    \item Si $X_3$ es casi irrelevante ($\beta_3 \approx 0$) o poco
      correlacionada con $X_2$ ($b_{32} \approx 0$), omitirla puede ser
      preferible.
    \item Si hay multicolinealidad severa ($r_{23}^2 \to 1$), el VIF infla la
      varianza del modelo largo y puede justificar aceptar el sesgo.
    \item En muestras grandes, el sesgo persiste pero la varianza tiende a
      cero; el estimador insesgado es casi siempre preferible.
    \item La teoría económica debe ser el primer criterio: si existe razón
      teórica sólida para incluir $X_3$, la insesgadez vale el costo de la
      varianza mayor.
  \end{enumerate}
\end{clave}

\vspace{14pt}
\noindent\rule{\linewidth}{0.4pt}

\noindent{\small\textit{Nota sobre la implementación en Stata.}
Dos restricciones del intérprete de Stata en archivos \texttt{.do} son
relevantes para este ejercicio: (i)~el operador~\texttt{;} no funciona como
separador de comandos fuera del modo \texttt{#delimit}; cada \texttt{local}
debe ocupar su propia línea. (ii)~La sintaxis \texttt{if \ldots\{\ \}\ else\ \{}
en una sola línea produce \texttt{r(198)} (\textit{code follows on the same line
as close brace}); la solución es reemplazar \texttt{if/else} por dos bloques
\texttt{if} independientes usando un escalar auxiliar de valor 0/1.
Adicionalmente, los programas definidos con \texttt{program define} no
persisten entre bloques de Org-Babel ejecutados en sesiones separadas, por lo
que el segundo bloque redefine \texttt{sim\_ej136} con \texttt{capture program drop}
antes de usarlo.}

\vspace{14pt}

\begin{center}
\small\textsc{Referencias}\\[6pt]
\end{center}

\noindent Gujarati, D.~N., \& Porter, D.~C. (2010). \textit{Econometría}
(5.a~ed.). McGraw-Hill. Capítulo~13, secciones~13.2--13.3, y apéndice~A.

\noindent Greene, W.~H. (2018). \textit{Econometric Analysis} (8.a~ed.).
Pearson. Capítulo~5.

\noindent Wooldridge, J.~M. (2016). \textit{Introductory Econometrics}
(6.a~ed.). Cengage. Capítulo~3.

\end{document}
